\documentclass[12pt,oneside]{article}
\usepackage[T1]{fontenc}
\usepackage{amsmath,amssymb,array,longtable}
\usepackage[english]{babel}
\usepackage[bookmarks=false,pdfauthor={Tomi Setsu},pdftitle={Tol: Syntax Specification}]{hyperref}
\usepackage{fancyhdr}

\setlength{\headheight}{15.2pt}
\pagestyle{fancy}
\fancyhf{}
\fancyfoot[C]{\thepage}
\renewcommand{\headrulewidth}{0.5pt}
\fancyhead[C]{Tol: Syntax Specification}

\title{Tol: Syntax Specification}
\author{Tomi Setsu}

\begin{document}
\maketitle

\begin{abstract}
This document describes the Tol surface syntax currently accepted by the
compiler in this repository. It is derived directly from
\verb|tol/lexer.cpp|, \verb|tol/ast-from-tokens.cpp|,
\verb|tol/pipe-lower-contract.cpp|, \verb|tol/pipeline.h|, and the bundled
standard library under \verb|crypto/smartcont/tol-stdlib/|. The target compiler
version is Tol \verb|1.3.0|.

Tol is TOS's own high-level smart-contract language: a typed, actor-oriented
extension of FunC for TVM contracts. It keeps the low-level TVM control that
FunC authors rely on, while adding automatic cell serialization, lazy loading,
typed data layouts, \verb|contract| blocks, typed \verb|receive(...)| dispatch,
deployment receivers, state machines, receiver-local query-id handling, and
Slice 3--6 domain standard libraries.

This text is descriptive of the implementation as it exists today. Where the
lexer reserves a token but the parser rejects or does not yet implement a
surface form, that is stated explicitly.
\end{abstract}

\section*{Introduction}
Tol is a statically typed source language for writing TOS smart contracts that
compile to TVM/Fift-style assembly. It is designed as a high-level extension
of FunC: the language keeps direct TVM/Fift-style escape hatches and explicit
control over cells, slices, builders, continuations, and registers, while
adding expressive typed data layouts, automatic serialization, lazy unpacking,
union types, nullable types, generic functions and structs, methods, and a
compiler-visible actor shape. Message dispatch, storage loading, state
transitions, getter methods, and several static-analysis checks are part of
the language surface.

At the parser level, a Tol source file is a sequence of version declarations,
imports, annotations, global declarations, constants, type aliases, functions,
structs, enums, and at most one entrypoint \verb|contract| declaration. The
compiler automatically loads \verb|@stdlib/common.tol| before the entrypoint
file.

\clearpage
\tableofcontents

\clearpage
\section{Overview}
The current Tol frontend has five relevant layers:

\begin{enumerate}
\item lexical analysis, implemented in \verb|tol/lexer.cpp|;
\item AST construction, implemented in \verb|tol/ast-from-tokens.cpp|;
\item symbol registration, type resolution, type inference, lvalue/rvalue
checking, pure/impure checking, and borrow checking;
\item TOS actor-layer checks and lowering, including state-machine,
field-scoping, receive-exhaustiveness, contract lowering, require-code
assignment, query-id propagation, reply correlation, and postponement checks;
\item code generation to Fif output.
\end{enumerate}

The compilation pipeline is ordered deliberately. In the current implementation
the contract checks run before \verb|pipeline_lower_contracts()|, while
\verb|pipeline_assign_require_codes()|,
\verb|pipeline_check_query_id_propagation()|,
\verb|pipeline_check_slice3_reply_correlation()|, and
\verb|pipeline_check_postponement()| run after contract lowering, when the
synthesized entrypoint bodies are visible.

\section{Relationship to FunC}
FunC exposes the TVM programming model directly: authors commonly write
entrypoint dispatch by hand, read an opcode from the inbound body, decode the
selected message, mutate persistent data in \verb|c4|, and construct outbound
actions explicitly. Tol is the higher-level TOS contract language built on that
model. It keeps the hand-written entrypoint style available:

\begin{verbatim}
struct (0x12345678) CounterIncrement {
    incBy: uint32
}

type AllowedMessage = CounterIncrement | CounterReset

fun onInternalMessage(in: InMessage) {
    val msg = lazy AllowedMessage.fromSlice(in.body);
    match (msg) {
        CounterIncrement => { ... }
        CounterReset => { ... }
        else => { ... }
    }
}
\end{verbatim}

A file may define \verb|fun onInternalMessage(in: InMessage)| directly,
provided it does not also declare a \verb|contract| block. The same actor shape
can be written as compiler-visible Tol contract syntax:

\begin{verbatim}
contract Counter {
    storage: CounterStorage

    @disclaim_query_id
    receive(msg: CounterIncrement) {
        save(CounterStorage { value: storage.value + msg.incBy });
    }
}
\end{verbatim}

The contract surface is more than a shorter spelling for manual dispatch. It
gives later compiler passes a closed set of receiver types and states, which
enables receive exhaustiveness, state reachability, field scoping, query-id
propagation, reply-correlation checks, and actor-layer hardening before code
generation.

\section{Lexical Structure}
\subsection{Whitespace and Comments}
Whitespace separates tokens and is otherwise ignored. The lexer recognizes:

\begin{itemize}
\item end-of-line comments beginning with \verb|//|;
\item block comments delimited by \verb|/*| and \verb|*/|.
\end{itemize}

Block comments are not nested. If a block comment reaches the end of the file
without a closing delimiter, the lexer treats the rest of the file as comment
text.

\subsection{Token Boundaries}
The lexer is implemented as a longest-prefix token trie. The following
single-character tokens are always token boundaries:

\begin{verbatim}
+  -  *  /  %  ?  :  ,  ;  (  )  [  ]  {  }  =
<  >  !  &  |  ^  ~  .
\end{verbatim}

The lexer also recognizes the multi-character tokens:

\begin{verbatim}
...  ==  !=  <=  >=  <=>  <<  >>  ~>>  ^>>
&&  ||  ~/  ^/
+=  -=  *=  /=  %=  &=  |=  ^=  <<=  >>=
->  =>  ++  --  ??
\end{verbatim}

The tokens \verb|++| and \verb|--| are recognized only to produce a targeted
compile-time error. Tol has no increment or decrement operator.

\subsection{Identifiers}
Ordinary identifiers begin with an ASCII letter, underscore, dollar sign, or a
backtick-delimited identifier. After the first character, ordinary identifiers
may contain ASCII letters, digits, underscores, or dollar signs:

\begin{verbatim}
counter
queryId
$generated
_privateName
\end{verbatim}

Backticks may wrap identifiers, including identifiers that would otherwise be
reserved words:

\begin{verbatim}
fun `do`(): int { return 0; }
var `with spaces` = 1;
\end{verbatim}

Backtick identifiers end at the next backtick before a newline. An unclosed
backtick is a lexical error.

\subsection{Integer Literals}
Integer literals are parsed as signed 257-bit values after tokenization. The
lexer accepts:

\begin{itemize}
\item decimal literals, for example \verb|0|, \verb|17|, \verb|1_000_000_000|;
\item hexadecimal literals with lowercase \verb|0x|, for example
\verb|0xff| and \verb|0x3b_9a_ca_00|;
\item binary literals with lowercase \verb|0b|, for example
\verb|0b0011_1011|.
\end{itemize}

Hexadecimal digits may be uppercase or lowercase, but the prefix itself is
lowercase. Underscore separators are allowed only between digits. A separator
immediately after \verb|0x| or \verb|0b| is rejected by the parser. Negative
literals are not separate lexical tokens; \verb|-1| is parsed as unary minus
applied to \verb|1| and then folded into a negative integer constant during
AST construction.

\subsection{String Literals}
Strings are delimited by ordinary double quotes or by triple double quotes.
Ordinary strings may not span lines. Triple-quoted strings may span lines and
are mainly used for long assembler bodies and codegen tests.

The parser unescapes:

\begin{verbatim}
\n  \r  \t  \\  \'  \"
\end{verbatim}

No FunC-style string suffixes are accepted. Compile-time string operations such
as CRC and SHA calculations are exposed through standard-library functions and
methods rather than through literal postfixes.

\subsection{Annotations}
An annotation token begins with \verb|@| and may contain alphanumeric
characters, underscores, and dots. The compiler gives special meaning to:

\begin{verbatim}
@pure  @inline  @inline_ref  @noinline  @method_id
@overflow1023_policy  @on_bounced_policy  @on
@deploy  @disclaim_query_id  @initial
@unknown_silent_drop  @unknown_throw
@implicit_protocol_default  @implicit_protocol_for
\end{verbatim}

The parser also allows opaque custom annotations of the forms
\verb|@custom|, \verb|@custom.*|, \verb|@test|, \verb|@test.*|, and
\verb|@deprecated|. Unknown non-custom annotations are rejected. Several of
the contract annotations above are accepted only in the specific contract
member positions described in Section~\ref{sec:contracts}.

\subsection{Reserved Keywords}
The lexer recognizes the following keywords:

\begin{verbatim}
fun  contract  receive  receive_external  storage
type  enum  struct  operator  infix
global  const  var  val  mutate  self
asm  builtin  private  readonly
true  false  null
return  repeat  do  while  break  continue
try  catch  throw  assert  if  else  match  lazy
as  is  tol  import  export
\end{verbatim}

Some keywords, such as \verb|operator|, \verb|infix|, and \verb|export|, are
reserved but currently rejected as unsupported top-level declarations.

\section{Types}
\label{sec:types}
\subsection{Type Expressions}
The parser uses the following grammar for type expressions:

\begin{verbatim}
TypeExpr ::= [ "|" ] NullableType { "|" NullableType } [ "->" TypeExpr ]

NullableType ::= SimpleType [ "<" TypeExpr { "," TypeExpr } [ "," ] ">" ]
                            [ "?" ]

SimpleType ::= identifier
             | "self"
             | "contract" | "receive" | "receive_external" | "storage"
             | "null"
             | "(" [ TypeExpr { "," TypeExpr } [ "," ] ] ")"
             | "[" [ TypeExpr { "," TypeExpr } [ "," ] ] "]"
\end{verbatim}

The leading \verb+|+ in a union is accepted, mirroring TypeScript-style union
notation. Function types are right-associative.

Examples:

\begin{verbatim}
int
int?
int | slice | null
Cell<MyStruct>
map<uint256, coins>
(int, slice)
[int, slice]
(int, slice) -> bool
int -> int -> int
\end{verbatim}

\subsection{Built-in and Standard Types}
The built-in type names are declared in \verb|@stdlib/common.tol|, which is
auto-imported. The main primitive and standard-library-backed types are:

\begin{verbatim}
int  bool  cell  slice  builder  continuation
address  any_address  never  unknown  void
intN  uintN  coins
varint16  varuint16  varint32  varuint32
string  bitsN  bytesN
map<K,V>  array<T>  lisp_list<T>  dict  tuple  Cell<T>
\end{verbatim}

Most primitive types occupy one TVM stack slot. Tensors occupy as many stack
slots as they have components. Structs are named tensor-like values at the
language level, but they may serialize into cells, builders, or slices through
generated pack/unpack code.

\subsection{Nullable Types and Null Safety}
The suffix \verb|?| forms a nullable type:

\begin{verbatim}
int?
Cell<Payload>?
address?
\end{verbatim}

The expression \verb|null| has the null type. The postfix \verb|!| asserts
non-nullness. Equality and inequality comparisons against \verb|null| are
rewritten into type checks internally, so \verb|x == null| and
\verb|x != null| participate in smart casts.

\subsection{Union Types and Smart Casts}
Union types are written with \verb+|+:

\begin{verbatim}
type AllowedMessage = Increment | Reset | null
\end{verbatim}

The operators \verb|is| and \verb|!is| test a value against a type and feed the
smart-cast analysis:

\begin{verbatim}
if (msg is Increment) {
    return msg.incBy;
}
\end{verbatim}

\subsection{Tensors, Arrays, and TVM Tuples}
Tol distinguishes:

\begin{itemize}
\item tensors, written \verb|(T1, T2, ...)|, which occupy several stack slots;
\item arrays, written \verb|array<T>| and commonly constructed with
\verb|[ ... ]|, which are backed by a TVM tuple value;
\item the raw \verb|tuple| type, an alias for \verb|array<unknown>|.
\end{itemize}

Tensor components and array-like shaped values may be accessed with dot-index
syntax:

\begin{verbatim}
var pair = (1, 2);
pair.1 = 3;
\end{verbatim}

\section{Top-Level Source Structure}
\label{sec:toplevel}
A source file is parsed as:

\begin{verbatim}
SourceItem ::= TolVersion
             | Import
             | Annotation
             | GlobalDecl
             | ConstDecl
             | TypeAlias
             | FunctionDecl
             | StructDecl
             | EnumDecl
             | ContractDecl
             | TopLevelGetter
             | ";"

Source ::= { SourceItem }
\end{verbatim}

Annotations are queued and then consumed by the following declaration. A
trailing annotation with no declaration is rejected.

\subsection{Version Declaration}
The version declaration is:

\begin{verbatim}
tol 1.3
tol 1.3.0
\end{verbatim}

Only exact version comparison is implemented. If the declared version differs
from the compiler version, the compiler emits a warning rather than rejecting
the source.

\subsection{Imports}
Imports have the form:

\begin{verbatim}
import "relative-file.tol"
import "@stdlib/dicts"
\end{verbatim}

Relative imports are resolved against the current file's directory. Imports
whose text starts with \verb|@| are resolved through configured path mappings;
\verb|@stdlib/...| is the standard-library mapping. Imported files are parsed
once per real path.

\subsection{Global Variables}
Global variable declarations have the form:

\begin{verbatim}
global name: Type
\end{verbatim}

Multiple global declarations in one statement are rejected. Assigning to a
global at declaration time is also rejected.

\subsection{Constants}
Constants have the form:

\begin{verbatim}
const NAME = expr
const NAME: Type = expr
\end{verbatim}

The initializer is parsed as an ordinary expression and later checked as a
constant expression by \verb|pipeline_check_constant_expressions()|. Multiple
constant declarations in one statement are rejected.

\subsection{Type Aliases}
Type aliases have the form:

\begin{verbatim}
type UserId = uint256
type Result<T, E = int> = Ok<T> | Err<E>
\end{verbatim}

The special form \verb|type map<K,V> = builtin| is used by the standard library
to declare compiler-backed types.

\section{Functions and Methods}
\label{sec:functions}
\subsection{Grammar}
The function grammar is:

\begin{verbatim}
FunctionDecl ::= "fun" [ ReceiverType "." ] name [ Generics ]
                 ParameterList [ ":" ReturnType ]
                 ( Block | AsmBody | "builtin" )

AsmBody ::= "asm" [ "(" { parameter-name } [ "->" { integer-literal } ] ")" ]
            string-literal { string-literal }

Generics ::= "<" GenericItem { "," GenericItem } [ "," ] ">"
GenericItem ::= identifier [ "=" TypeExpr ]

ParameterList ::= "(" [ Parameter { "," Parameter } [ "," ] ] ")"
Parameter ::= [ "mutate" ] (identifier | "self" | "_")
              [ ":" TypeExpr ] [ "=" Expr ]
\end{verbatim}

Parameter types are mandatory for ordinary functions except for the special
\verb|self| parameter of methods. Lambda parameters may omit types.
\verb|mutate| parameters cannot have default values.

If a return type is omitted, the compiler infers it. Omitted return type does
not mean \verb|void|.

\subsection{Methods}
A function whose name is preceded by a receiver type and a dot is a method:

\begin{verbatim}
fun Point.sum(self): int {
    return self.x + self.y;
}

fun builder.storePair(mutate self, x: int, y: int): void {
    self.storeInt(x, 32);
    self.storeInt(y, 32);
}
\end{verbatim}

The first parameter may be \verb|self|. Without \verb|mutate|, \verb|self| is
immutable. A method may declare \verb|: self| as its return type only when it
accepts \verb|self| as its first parameter.

\subsection{Entrypoint Names}
The following top-level function names are reserved entrypoints when they have
no receiver:

\begin{verbatim}
main  onInternalMessage  onExternalMessage  onRunTickTock
onSplitPrepare  onSplitInstall  onBouncedMessage
\end{verbatim}

FunC/Fift-style names such as \verb|recv_internal| and \verb|recv_external|
are rejected with a diagnostic telling the author to use
\verb|onInternalMessage|.

\subsection{Assembler and Builtin Bodies}
An assembler function has the form:

\begin{verbatim}
fun f(x: int): int asm "INC";
fun g(a: int, b: int): (int, int) asm(a b -> 1 0) "SWAP";
\end{verbatim}

An \verb|asm| function must specify a return type. The optional order clause
lists argument names and returned stack indexes. One or more string literals
form the assembler command body. The parser rejects assembler functions with
more than 16 formal parameters.

The \verb|builtin| body form is reserved for compiler-backed functions in the
standard library:

\begin{verbatim}
fun slice.loadInt(mutate self, len: int): int
    builtin
\end{verbatim}

\subsection{Annotations on Functions}
The function annotations accepted by the parser are:

\begin{itemize}
\item \verb|@pure|;
\item \verb|@inline|;
\item \verb|@inline_ref|;
\item \verb|@noinline|;
\item \verb|@method_id(expr)| on regular top-level functions, where
\verb|expr| must resolve to an integer method id or a string method name;
\item \verb|@on_bounced_policy("manual")| or
\verb|@on_bounced_policy("ignore")| on \verb|onInternalMessage|.
\end{itemize}

Opaque custom annotations are allowed but ignored by the compiler.

\subsection{File-Scope Get Methods}
At top level, \verb|get fun name(...)| declares a legacy file-scope get method:

\begin{verbatim}
get fun seqno(): uint32 {
    return 0;
}
\end{verbatim}

File-scope getters must be ordinary code-block functions: no receiver type, no
generic parameters, no \verb|mutate| parameters, and no \verb|asm| or
\verb|builtin| body. A file containing a \verb|contract| block cannot also
declare a file-scope \verb|get fun|.

\section{Structs, Enums, and Serialization Prefixes}
\label{sec:structs}
\subsection{Struct Declarations}
Struct declarations have the form:

\begin{verbatim}
struct Name [ Generics ] [ StructBody ]
struct (PrefixExpr) Name [ Generics ] StructBody

StructBody ::= "{" { StructField [ "," | ";" ] } "}"
StructField ::= [ FieldAnnotations ] [ "private" ] [ "readonly" ]
                name ":" TypeExpr [ FieldAnnotations ] [ "=" Expr ]
\end{verbatim}

If the body is omitted, the struct is empty.

A struct declaration may carry \verb|@overflow1023_policy("suppress")|. This
suppresses the static warning that auto-serialization can theoretically exceed
the TVM cell limits of 1023 bits or four references. The annotation does not
change the wire layout; an actually oversized value can still fail at runtime.

\subsection{Serialization Prefixes}
The optional expression in \verb|struct (PrefixExpr) Name| declares a
serialization prefix. For ordinary TOS messages, a 32-bit prefix is used as the
message opcode:

\begin{verbatim}
const OP_TRANSFER: uint32 = 0x0f8a7ea5

struct (OP_TRANSFER) Transfer {
    queryId: uint64
    amount: coins
}
\end{verbatim}

Tol permits prefixes that are not 32 bits for general serialization and union
encoding. TOS \verb|contract receive(...)| declarations are stricter: every
received message type must resolve to exactly a 32-bit prefix.

\subsection{Struct Field Modifiers}
Fields may be \verb|private| and \verb|readonly|:

\begin{verbatim}
struct Cursor {
    private readonly t: tuple
    index: uint16 = 0
}
\end{verbatim}

\verb|private| fields are accessible only through methods on that struct.
\verb|readonly| fields cannot be modified after construction.

State-bearing contracts may also use \verb|@on(State1, State2)| to restrict a
storage field to a set of states:

\begin{verbatim}
struct AuctionStorage {
    winner: address @on(Closed)
}
\end{verbatim}

The \verb|@on| annotation may appear before the field or after its type. It is
valid only on fields of the storage struct of a state-bearing contract. Field
default expressions in non-generic structs are checked as compile-time
constant expressions.

\subsection{Object Literals and Spread}
Object literals are written either with an explicit type or with an inferred
type:

\begin{verbatim}
val p1: Point = { x: 1, y: 2 };
val p2 = Point { x: 1, y: 2 };
val p3 = { ...p2, y: 3 };
\end{verbatim}

The shorthand field syntax \verb|{ x, y }| expands to
\verb|{ x: x, y: y }|.

\subsection{Enums}
Enum declarations have the form:

\begin{verbatim}
enum Color {
    Red,
    Green,
    Blue = 7,
}

enum ErrorClass: uint8 {
    Ok = 0,
    Protocol = 4,
}
\end{verbatim}

Members may omit explicit values. The optional colon type fixes the serialized
integer type and must resolve to an integer-width type such as \verb|uint8| or
to \verb|coins|.

\section{Statements}
\label{sec:statements}
\subsection{Grammar}
Function and receiver bodies contain:

\begin{verbatim}
Stmt ::= LocalVarDecl
       | Block
       | "return" [ Expr ]
       | "if" "(" Expr ")" Block [ "else" (IfStmt | Block) ]
       | "repeat" "(" Expr ")" Block
       | "while" "(" Expr ")" Block
       | "do" Block "while" "(" Expr ")"
       | "throw" Expr
       | "throw" "(" Expr [ "," Expr ] ")"
       | "assert" "(" Expr "," Expr ")"
       | "assert" "(" Expr ")" "throw" Expr
       | "try" Block "catch" [ "(" CatchVars ")" ] Block
       | Expr
       | ";"
\end{verbatim}

Inside contract receivers, the pseudo-statements \verb|become State| and
\verb|keep_state| are also recognized.

Within a block, statements that do not end with a braced statement require a
semicolon separator. A final statement before \verb|}| may omit the semicolon
when the parser can see the closing brace.

\subsection{Local Variables}
Local variables use \verb|var| or \verb|val|:

\begin{verbatim}
var x = 1;
val y: uint32 = 2;
var (a, b) = (1, 2);
var [first, second] = [1, 2];
var late: cell;
\end{verbatim}

\verb|val| bindings are immutable. A local declaration without an initializer
is accepted only in the late-init form with an explicit type.

\subsection{Throw, Assert, and Try/Catch}
\verb|throw| accepts either an exit code or an exit code plus argument:

\begin{verbatim}
throw 0xffff;
throw (0xffff, payload);
\end{verbatim}

\verb|assert| accepts either a comma form or a FunC-like \verb|throw| form:

\begin{verbatim}
assert(x > 0, 100);
assert(x > 0) throw 100;
\end{verbatim}

A \verb|catch| clause may bind no variables, one variable, or two variables:

\begin{verbatim}
try { ... } catch { ... }
try { ... } catch (excNo) { ... }
try { ... } catch (excNo, arg) { ... }
\end{verbatim}

\subsection{Unsupported Loop Control}
The lexer reserves \verb|break| and \verb|continue|, but the parser currently
rejects both with a diagnostic. Loop bodies must be structured without them.

\section{Expressions}
\label{sec:expressions}
\subsection{Primary Forms}
The highest-precedence expression forms are:

\begin{verbatim}
Primary ::= "()"
          | "(" Expr ")"
          | "(" Expr "," Expr { "," Expr } [ "," ] ")"
          | "[" [ Expr { "," Expr } [ "," ] ] "]"
          | integer-literal
          | string-literal
          | "_"
          | "true" | "false" | "null" | "self"
          | identifier [ InstantiationTypes ]
          | TypeName ObjectBody
          | TypeName "[" ... "]"
          | ObjectBody
          | MatchExpr
          | LambdaExpr
          | LazyExpr
\end{verbatim}

The empty parenthesized expression \verb|()| is the unit tensor. Square
brackets construct an array-like value backed by a TVM tuple. An object body
without an explicit type relies on contextual type inference.

\subsection{Calls, Methods, and Postfix Operators}
Function calls and method calls are left-associative:

\begin{verbatim}
f(1, 2)
builder.storeUint(1, 32).storeUint(2, 32)
genericFunction<int>(x)
obj.method<uint32>(arg)
\end{verbatim}

Arguments may be marked \verb|mutate| when the callee declares the
corresponding parameter \verb|mutate|:

\begin{verbatim}
state.receiveBid(mutate queue, bid);
\end{verbatim}

The postfix non-null assertion is \verb|!|:

\begin{verbatim}
value!.field
\end{verbatim}

\subsection{Casts and Type Tests}
The cast operator is \verb|as|:

\begin{verbatim}
val small = x as uint32;
\end{verbatim}

Type tests use \verb|is| and \verb|!is|:

\begin{verbatim}
if (msg is Transfer) { ... }
if (msg !is Transfer) { ... }
\end{verbatim}

\subsection{Match Expressions}
\verb|match| is an expression:

\begin{verbatim}
val result = match (subject) {
    Transfer => subject.amount,
    Reset => 0,
    else => throw 0xffff,
};
\end{verbatim}

Each arm pattern is parsed either as a type, as a constant expression, or as
\verb|else|. Unbraced expression arms require comma separation; braced arms may
omit the comma.

The subject may also be a local declaration:

\begin{verbatim}
match (var msg = lazy AllowedMessage.fromSlice(in.body)) {
    Transfer => { ... }
    else => { ... }
}
\end{verbatim}

\subsection{Lambdas and Lazy Loading}
Lambda expressions have the form:

\begin{verbatim}
fun(x: int): int { return x + 1; }
\end{verbatim}

Parameter types may be omitted in lambdas when inference can supply them.

The \verb|lazy| operator delays deserialization and field loading:

\begin{verbatim}
val msg = lazy Transfer.fromSlice(in.body);
val storage = lazy Storage.fromCell(contract.getData());
\end{verbatim}

\subsection{Operator Precedence}
The parser implements the following precedence levels, from highest to lowest:

\begin{center}
\begin{tabular}{|l|m{0.68\textwidth}|}
\hline
level & forms \\ \hline
100 & primary expressions \\ \hline
90 & function call and postfix \verb|!| \\ \hline
80 & dot access, method call, and numeric component access \\ \hline
75 & unary \verb|!|, \verb|~|, \verb|-|, \verb|+| \\ \hline
40 & \verb|as|, \verb|is|, \verb|!is| \\ \hline
30 & \verb|*|, \verb|/|, \verb|%|, \verb|~/|, \verb|^/| \\ \hline
20 & \verb|+|, \verb|-| \\ \hline
17 & \verb|<<|, \verb|>>|, \verb|~>>|, \verb|^>>| \\ \hline
15 & one comparison: \verb|==|, \verb|!=|, \verb|<|, \verb|>|,
\verb|<=|, \verb|>=|, \verb|<=>| \\ \hline
14 & bitwise \verb|&|, \verb+|+, \verb|^| \\ \hline
13 & logical \verb|&&|, \verb+||+ \\ \hline
10 & assignment, compound assignment, ternary \verb|?:|, null coalescing
\verb|??| \\ \hline
\end{tabular}
\end{center}

Comparison operators are not chainable. Assignment, ternary, and null
coalescing are parsed right-to-left. The parser emits strict diagnostics for
several ambiguous precedence patterns, such as mixing bitwise operators with
comparisons without parentheses.

\section{Contract Blocks}
\label{sec:contracts}
\subsection{Purpose}
The TOS actor-layer extension promotes the common message-dispatch pattern into
compiler-visible syntax:

\begin{verbatim}
contract Name {
    storage: StorageType

    receive(msg: MessageType) {
        ...
    }

    get fun getter(): uint64 {
        ...
    }
}
\end{verbatim}

The lowering synthesizes \verb|loadData|, \verb|save|, and
\verb|onInternalMessage|. If the contract declares external receivers, it also
synthesizes \verb|onExternalMessage|.

\subsection{Contract Grammar}
The accepted member forms are:

\begin{verbatim}
ContractDecl ::= [ ContractAnnotations ] "contract" name "{"
                 ContractMember*
                 "}"

ContractMember ::= "storage" ":" StorageType [ ";" ]
                 | "states" ":" State { "," State } [ ";" ]
                 | "@initial" "state" State [ ";" ]
                 | "@unknown_silent_drop" [ ";" ]
                 | "@unknown_throw" "(" integer-literal ")" [ ";" ]
                 | "@implicit_protocol_default" [ ";" ]
                 | "@implicit_protocol_for" "(" MessageType "," State ")" [ ";" ]
                 | [ "@deploy" ] [ "@disclaim_query_id" ] ReceiveBlock
                 | ReceiveExternalBlock
                 | [ GetterAnnotations ] "get" "fun" ...
\end{verbatim}

The only contract annotation accepted before the \verb|contract| keyword is
\verb|@on_bounced_policy("manual")| or \verb|@on_bounced_policy("ignore")|.
\verb|StorageType| must be a simple type identifier, not a generic, union,
nullable, or shaped type expression.

A contract block must contain exactly one \verb|storage:| declaration and at
least one internal \verb|receive(...)| declaration. The storage declaration
must appear before receivers and getters. Free-standing \verb|fun|
declarations are not permitted inside a contract block.

At the file level, the current parser supports one \verb|contract| declaration
per \verb|.tol| file. A file containing a contract cannot also contain a
top-level \verb|onInternalMessage|. If the contract declares
\verb|receive_external|, the file cannot also contain a top-level
\verb|onExternalMessage|.

\subsection{Typed Internal Receivers}
Internal receivers have the form:

\begin{verbatim}
receive(msg: Transfer) {
    ...
}

receive(msg: Transfer) on Open {
    ...
}
\end{verbatim}

The message type must be a simple struct type with exactly a 32-bit
serialization prefix. The parameter name is bound to a lazy decode of
\verb|in.body|. Normal receivers also have \verb|storage| and \verb|in| in
scope.

The lowering dispatches on the first 32 bits of the body. It does not pre-load
\verb|queryId| as a global dispatch field; query-id semantics are receiver
local.

Internal receiver opcodes must be unique. In a state-bearing contract, the
same message struct and opcode may be used by multiple receivers only when
each receiver has a distinct \verb|on State| clause. A duplicate
opcode/state pair, or a mix of state-qualified and unqualified receivers for
the same opcode, is rejected.

\subsection{Deployment Receivers}
\verb|@deploy| marks the bootstrap receiver:

\begin{verbatim}
@deploy
receive(msg: Deploy) {
    save(Storage { owner: msg.owner });
}
\end{verbatim}

The deploy branch runs before \verb|loadData()| so that an empty c4 can be
initialized. Inside a deploy receiver, \verb|storage| is not in scope. The
receiver must materialize storage by calling \verb|save(...)|.

\verb|@deploy| is mutually exclusive with an \verb|on State| clause and with
\verb|@disclaim_query_id|. At most one deploy receiver is allowed. In a
state-bearing contract, a deploy receiver requires an \verb|@initial state|
declaration so that the compiler can inject the initial hidden state tag into
the saved storage value.

\subsection{Unknown Opcode Handling}
There are four internal-message unknown-opcode modes:

\begin{itemize}
\item default mode: throw \verb|OP_ERROR|, currently \verb|0x00010001|;
\item \verb|@unknown_silent_drop;|: return normally without actions;
\item \verb|@unknown_throw(N);|: throw the specified code;
\item \verb|receive(msg: UnknownOpcode) { ... }|: run a catch-all body with
the raw \verb|in.body| slice bound to the parameter.
\end{itemize}

The catch-all pseudo-type \verb|UnknownOpcode| has no serialization encoding
and cannot be used as an ordinary message body type. It is mutually exclusive
with \verb|@unknown_silent_drop| and \verb|@unknown_throw|. If no visible
unknown-opcode policy is declared, the compiler keeps the default Protocol
throw but emits a receive-exhaustiveness warning. These unknown-opcode
annotations apply to the internal-message dispatch domain only.

\subsection{External Receivers}
External receivers have the form:

\begin{verbatim}
receive_external(msg: SignedRequest) {
    ...
}
\end{verbatim}

They synthesize \verb|onExternalMessage| and dispatch on a 32-bit body prefix.
Normal external receivers bind the parameter to a lazy decode of the external
body and load \verb|storage| before running the user body. They do not have
internal-message \verb|InMessage| context and do not participate in query-id
propagation checks. State-machine guards apply only to internal receivers.

\verb|receive_external(msg: UnknownOpcode) { ... }| is accepted as an
external-only catch-all; its parameter is the raw \verb|inMsgBody| slice. If no
external catch-all is declared, an unrecognized external opcode throws the
fixed Protocol sentinel \verb!(4 << 16) | 0xffff!. The external catch-all does
not synthesize the normal \verb|storage| local. External opcodes must be unique
within the external dispatch table, but they may collide with internal receiver
opcodes because the entrypoints are separate.

\subsection{Get Methods}
Contract-local getters use:

\begin{verbatim}
get fun getBalance(): coins {
    return storage.balance;
}

@method_id(0x1001)
get fun getOwner(): address {
    return storage.owner;
}
\end{verbatim}

A \verb|get fun| must declare an explicit return type. The compiler derives the
method id from the getter name unless \verb|@method_id(expr)| pins it.
Method-id collisions are checked per contract.

Contract-local getters accept \verb|@method_id|, \verb|@inline|,
\verb|@inline_ref|, \verb|@noinline|, \verb|@pure|, and opaque custom
annotations such as \verb|@deprecated|. Getter bodies are read-only: calls that
commit storage or emit actions, such as \verb|save(...)|,
\verb|contract.setData(...)|, \verb|sendRawMessage(...)|, \verb|.send(...)|,
or \verb|.sendAndEstimateFee(...)|, are rejected. Getter names may not shadow
generated contract internals such as \verb|loadData|, \verb|save|,
\verb|onInternalMessage|, or \verb|onExternalMessage|.

\section{State Machines}
\label{sec:state-machines}
\subsection{State Declaration}
A state-bearing contract declares:

\begin{verbatim}
contract Auction {
    storage: AuctionStorage
    states: Open, Closed
    @initial state Open

    receive(msg: Bid) on Open {
        ...
        keep_state;
    }

    receive(msg: Close) on Open {
        ...
        become Closed;
    }
}
\end{verbatim}

The compiler injects a hidden state tag into the serialized storage. User code
may not declare or access a field named \verb|__state|. A state-bearing
contract must declare exactly one \verb|@initial state|, and every ordinary
internal receiver must declare an \verb|on State| clause.

\subsection{Transitions}
Every control-flow path in a state-bearing receiver must end in either:

\begin{itemize}
\item \verb|become State|, which writes a new hidden state tag;
\item \verb|keep_state|, which preserves the current state tag.
\end{itemize}

Both forms are tail-position statements. Code after either form is rejected as
unreachable or invalid. \verb|become| accepts only a static state identifier.
\verb|become| and \verb|keep_state| are not valid inside \verb|@deploy|
receivers; deploy receivers use \verb|@initial state| as the authoritative
initial tag. They are also rejected inside \verb|receive(msg: UnknownOpcode)|
because the catch-all body has no source state.

\subsection{State Checks}
The compiler performs:

\begin{itemize}
\item state reachability checks from the initial state;
\item receive exhaustiveness warnings over the declared message-by-state matrix;
\item field-scoping checks for \verb|@on(...)| storage fields;
\item hidden state access checks.
\end{itemize}

An intentionally missing known-opcode/wrong-state pair may be documented with:

\begin{verbatim}
@implicit_protocol_for(MessageType, StateName);
\end{verbatim}

or, for every missing pair:

\begin{verbatim}
@implicit_protocol_default;
\end{verbatim}

\section{Actor-Layer Static Analysis}
\label{sec:actor-static-analysis}
\subsection{Query-Id Propagation}
TOS message policy treats \verb|queryId| as the request/reply correlation
field for envelope-shaped bodies. The compiler checks receiver scopes
independently after contract lowering. A fire-and-forget receiver may state its
intent explicitly:

\begin{verbatim}
@disclaim_query_id
receive(msg: BroadcastEvent) {
    ...
}
\end{verbatim}

One receiver's disclaimer does not silence warnings in a sibling receiver.
The lower-level \verb|disclaim_query_id()| builtin remains available in
hand-written \verb|onInternalMessage| code.

\subsection{Require Code Assignment}
\verb|require(cond, ErrorClass.X)| is a standard-library/compiler surface used
by receiver bodies. After contract lowering, the compiler assigns stable
per-contract error codes to require sites without explicit third arguments and
emits a require-site manifest comment in generated Fif output. Authors may pin
a code explicitly:

\begin{verbatim}
require(amount <= storage.balance, ErrorClass.Permanent, 0x1201);
\end{verbatim}

\subsection{Storage and c4 Escape Restrictions}
Inside \verb|contract| receiver bodies, the compiler rejects several raw c4
escape paths that would bypass the declared storage and state model. Examples
include direct \verb|contract.getData()| access, wrappers that reach c4, and
assembler bodies that mention c4/control-register aliases. Authors who need a
raw migration path should use the lower-level hand-written entrypoint style
instead of a \verb|contract| block.

\subsection{Postponement and Slice Checks}
The Slice 4 postponement checker rejects unsafe postponement paths, including
external-message enqueue attempts and direct queue internals that bypass the
bounded stdlib helper. Slice 3 and Slice 5 checks cover reply-correlation and
domain-pattern hardening surfaces used by the bundled standard library.

\section{Standard Library Surface}
\label{sec:stdlib}
\subsection{Common Prelude}
\verb|@stdlib/common.tol| is loaded before the entrypoint file. It declares
the built-in types, core TVM wrappers, arithmetic and logical helpers,
builders, slices, cells, address helpers, message creation, send modes, reserve
modes, random helpers, debugging helpers, and contract/blockchain namespaces.

The standard library presents many low-level FunC and TVM operations as typed
methods:

\begin{verbatim}
cell.hash()
slice.loadUint(32)
builder.storeUint(value, 32)
contract.getData()
blockchain.now()
\end{verbatim}

\subsection{Explicit Imports}
More specialized packages require explicit imports. The current repository
contains, among others:

\begin{verbatim}
@stdlib/tvm-dicts
@stdlib/tvm-lowlevel
@stdlib/strings
@stdlib/gas-payments
@stdlib/reflection
@stdlib/exotic-cells
@stdlib/lisp-lists
@stdlib/slice3-common
@stdlib/jetton
@stdlib/nft
@stdlib/wallet
@stdlib/multisig
@stdlib/ownable
@stdlib/postponement
@stdlib/auction
@stdlib/governance
@stdlib/oracle
@stdlib/payment-channel
@stdlib/delivery
@stdlib/time
@stdlib/schedule
@stdlib/supervision
@stdlib/capability
@stdlib/safe-payments
\end{verbatim}

Imports are explicit; a test or contract that uses symbols from a package
directly must import that package, even if an imported production source also
imports it.

\subsection{Slice-Layer Packages}
The TOS-specific stdlib is staged:

\begin{itemize}
\item Slice 1: envelope, error class, \verb|OP_ERROR|, extra-flags constants,
and query-id discipline;
\item Slice 3: first-wave patterns for jetton, NFT, wallet, multisig, and
ownable contracts;
\item Slice 4: bounded postponement as explicit contract-owned storage;
\item Slice 5: auction, governance, oracle, and payment-channel patterns plus
ABI manifest validation;
\item Slice 6: delivery records, scheduled-message shapes, time helpers,
monitor/supervision helpers, capability grants, and safe-payment helpers.
\end{itemize}

Some Slice 6 protocol behavior remains activation-gated even though the
repo-side Tol and stdlib surfaces are present.

\section{Reserved but Unsupported Surface}
The following tokens or forms are reserved or lexed but currently rejected:

\begin{itemize}
\item \verb|export|, \verb|operator|, and \verb|infix| as top-level
declarations;
\item \verb|break| and \verb|continue|;
\item \verb|++| and \verb|--|;
\item multiple \verb|contract| declarations in one \verb|.tol| file;
\item a \verb|contract| block mixed with top-level \verb|onInternalMessage|;
\item free-standing \verb|fun| declarations inside a contract block;
\item receiver message prefixes that are not exactly 32 bits;
\item \verb|receive_external(...) on State|;
\item \verb|@deploy| combined with \verb|on State| or
\verb|@disclaim_query_id|.
\end{itemize}

\appendix
\section{Summary Grammar}
This appendix collects the main grammar fragments in one place.

\begin{verbatim}
Source ::= { SourceItem }

SourceItem ::= TolVersion
             | Import
             | Annotation
             | GlobalDecl
             | ConstDecl
             | TypeAlias
             | FunctionDecl
             | StructDecl
             | EnumDecl
             | ContractDecl
             | TopLevelGetter
             | ";"

TolVersion ::= "tol" semver
Import     ::= "import" string-literal

GlobalDecl ::= "global" identifier ":" TypeExpr
ConstDecl  ::= "const" identifier [ ":" TypeExpr ] "=" Expr
TypeAlias  ::= "type" identifier [ Generics ] "=" (TypeExpr | "builtin")

FunctionDecl ::= "fun" [ TypeExpr "." ] identifier [ Generics ]
                 ParameterList [ ":" (TypeExpr | "self") ]
                 ( Block | AsmBody | "builtin" )

AsmBody ::= "asm" [ "(" { identifier | "self" }
                    [ "->" { integer-literal } ] ")" ]
            string-literal { string-literal }

TopLevelGetter ::= "get" "fun" identifier ParameterList
                   [ ":" TypeExpr ] Block

Generics ::= "<" GenericItem { "," GenericItem } [ "," ] ">"
GenericItem ::= identifier [ "=" TypeExpr ]

ParameterList ::= "(" [ Parameter { "," Parameter } [ "," ] ] ")"
Parameter ::= [ "mutate" ] (identifier | "self" | "_")
              [ ":" TypeExpr ] [ "=" Expr ]

StructDecl ::= "struct" [ "(" Expr ")" ] identifier [ Generics ]
               [ "{" { StructField [ "," | ";" ] } "}" ]

StructField ::= [ Annotation* ] [ "private" ] [ "readonly" ]
                identifier ":" TypeExpr [ Annotation* ] [ "=" Expr ]

EnumDecl ::= "enum" identifier [ ":" TypeExpr ]
             "{" { EnumMember [ "," | ";" ] } "}"
EnumMember ::= identifier [ "=" Expr ]

ContractDecl ::= [ "@on_bounced_policy" "(" string-literal ")" ]
                 "contract" identifier "{" { ContractMember } "}"

ContractMember ::= "storage" ":" identifier [ ";" ]
                 | "states" ":" identifier { "," identifier } [ ";" ]
                 | "@initial" "state" identifier [ ";" ]
                 | "@unknown_silent_drop" [ ";" ]
                 | "@unknown_throw" "(" integer-literal ")" [ ";" ]
                 | "@implicit_protocol_default" [ ";" ]
                 | "@implicit_protocol_for" "(" identifier "," identifier ")" [ ";" ]
                 | [ "@deploy" ] [ "@disclaim_query_id" ] ReceiveBlock
                 | ReceiveExternalBlock
                 | [ Annotation* ] "get" "fun" identifier ParameterList
                   ":" TypeExpr Block

ReceiveBlock ::= "receive" "(" identifier ":" identifier ")"
                 [ "on" identifier ] Block

ReceiveExternalBlock ::= "receive_external"
                         "(" identifier ":" identifier ")" Block

TypeExpr ::= [ "|" ] NullableType { "|" NullableType } [ "->" TypeExpr ]
NullableType ::= SimpleType [ "<" TypeExpr { "," TypeExpr } [ "," ] ">" ]
                            [ "?" ]
SimpleType ::= identifier
             | "self"
             | "contract" | "receive" | "receive_external" | "storage"
             | "null"
             | "(" [ TypeExpr { "," TypeExpr } [ "," ] ] ")"
             | "[" [ TypeExpr { "," TypeExpr } [ "," ] ] "]"
\end{verbatim}

\section{Expression Precedence Summary}
\begin{verbatim}
100  primary expressions
 90  function call, postfix !
 80  dot access, method call, numeric component access
 75  unary !  ~  -  +
 40  as  is  !is
 30  *  /  %  ~/  ^/
 20  +  -
 17  <<  >>  ~>>  ^>>
 15  ==  !=  <  >  <=  >=  <=>
 14  &  |  ^
 13  &&  ||
 10  =  +=  -=  *=  /=  %=  <<=  >>=  &=  |=  ^=  ?:  ??
\end{verbatim}

\section{Reference Sources}
The implementation facts in this document are taken from the local repository:

\begin{itemize}
\item \verb|tol/lexer.cpp|
\item \verb|tol/ast-from-tokens.cpp|
\item \verb|tol/pipe-lower-contract.cpp|
\item \verb|tol/tol.cpp|
\item \verb|tol/pipeline.h|
\item \verb|crypto/smartcont/tol-stdlib/common.tol|
\item \verb|crypto/smartcont/tol-stdlib/*.tol|
\item \verb|doc/tos-language-syntax-policy.md|
\end{itemize}

\end{document}
